\chapter{Limitations and open points}
\label{ch:limits}

The items below are properties of the delivered RTL, not defects. They are listed
so that integration proceeds with them known.

\section{Instruction set}

Multiply and divide are not implemented; the \sig{M} extension was out of scope
per the specification and its encodings decode as illegal instructions. There are
no atomic instructions --- RV32I has no load-reserved / store-conditional and none
were added. There are no floating-point, compressed or vector extensions.
\reg{misa} reports exactly what is built.

Only machine mode exists. \reg{mstatus.MPP} is hardwired to \texttt{11} because
there is no other privilege level to return to. There is no virtual memory and no
MMU or PMP.

\section{Memory system}

There are no caches, no write buffer and no coherence mechanism. Every load and
store goes to the bus and blocks the execute stage for the whole round trip.
Consequently a hart is sequentially consistent by construction, which is what
makes flag-based sharing between cores sound, but single-thread memory
performance is bounded by bus latency.

Each master port carries at most one outstanding transaction. AXI4-Lite has no
bursts or IDs, and none are used. The per-channel handshake tracking in the LSU
was written so that a later move to AXI4-Full is additive, but that move has not
been made.

\section{Peripheral address decoding}

The peripheral slaves decode a byte offset of eight bits and no more. The system
map resolves this with exact 256\,byte windows, so anything above a block's
register file misses every window and is answered \sig{DECERR}; the alternative,
a wider comparison inside each shared slave, would have meant editing blocks
owned by other people. What remains is that a block reached through some other
interconnect, with a window wider than its register file, would alias silently.
The window size is a property of the map, not of the block.

\section{Interrupts}

The PIC serves one CPU. Source-to-core targeting for a multi-core configuration
is undefined, and each core would want its own timer compare register.

The PIC supports preemptive nesting natively, and the feature bench exercises it.
Reaching a second level through the CPU needs handler software that re-enables
\reg{MIE} after saving \reg{mepc}, \reg{mcause} and \reg{mstatus} --- the standard
RISC-V sequence --- because trap entry clears \reg{MIE}.

This was described as a software addition on top of an end-of-interrupt path
that already worked. It was not: the core tracked one in-progress bit, so
\sig{MRET} could not tell an interrupt return from an exception return. Two
nested handlers produced one end-of-interrupt instead of two and left the outer
level on the PIC's stack forever, and an exception taken inside a handler sent
its end-of-interrupt early. The core now carries one bit per open trap level and
\file{tb\_traps} covers both cases, so the remaining part really is software.

The depth of that stack is 16, which is the largest \reg{NEST\_MAX} the PIC
accepts, so no reachable interrupt nesting can overflow it. Nothing bounds how
deep a chain of synchronous exceptions inside handlers can go; past 16 open
levels the outermost level's kind is lost.

\section{Reset}

Reset is asynchronous and active low across all three blocks. Release is not
synchronised. In simulation this is harmless, and the benches deliberately
release reset between clock edges so the behaviour is exercised rather than
avoided. Real silicon would require a reset synchroniser --- asynchronous
assertion, synchronous de-assertion --- to respect recovery and removal timing.

Reset asserted in the middle of a bus transaction is not exercised: the models
and the CPU reset together. This becomes relevant once the system defines its
reset controller.

The dual-port SRAM in the wider system uses synchronous reset while the CPU and
DMA use asynchronous reset. Harmonising that is a system-level decision.

\section{Counters and CSR conformance}

\reg{mcountinhibit} is not implemented, which the specification permits; software
therefore cannot freeze a counter around a measurement window.
\reg{mtvec.MODE} does not force a legal value on a reserved write, as noted in
Chapter~\ref{ch:csr}. The performance-counter events are fixed in hardware and
\reg{mhpmevent} is tied off, so the events cannot be reprogrammed.

\section{Verification scope}

Verification is directed and self-checking: expected values are written by
hand, so the results establish that the delivery behaves as this document
describes, at the values listed in Section~\ref{sec:valscope}. Functional
coverage stands at 88 of 92 bins, with all four misses accounted for
individually in Table~\ref{tab:covmiss}. Section~\ref{sec:valscope} states, area
by area, which value spaces were covered exhaustively and which were sampled;
the rows marked \emph{selected} are where the delivery is thinnest. The
testbench interconnect covers two slaves; richer fabrics, including one that
generates its own decode errors, arrive with the real interconnect.

\section{Area, timing and power}

No synthesis has been run, so this document contains no area, frequency or power
figure for any block. One design choice is known to carry a cost that only
synthesis can quantify: the single decode-and-execute stage produces a long
combinational path (Chapter~\ref{ch:cpuarch}), which would set the achievable
clock frequency on a real target.
